% !TeX root = ../papers/nc-p3.tex
\section{特異な楕円核の非分岐鉛筆での二十台点の排除}
\label{sec:singular-residual-pencils}
\begin{theorem}\label{thm:singular-residual-pencils}
四辺有限被約な標準有限Frobenius窓で、整な八次曲面の八余因子核が節点または尖点を持つとする。
前節の有限非分岐元は特殊な $T$ を含めて全て不可能である。
従って節点では二捩れまたは境界、尖点では境界Hankel行列だけが残る。
\end{theorem}
\begin{proof}[正規化上の三つの共通根]
前節の二変数式で二項四次式 $a,b$ の最大公約因子を $g$ とし、
残余を互いに素な次数 $m=4-\deg g$ の二項式 $A,B$ とする。
$g(z)g(w)$ を除いた核行列の階数は $m$ 以上である。
節点では必要なら $A,B$ を入れ替え（$t$ も逆にして全体倍率を取る）、
$A$ の最高係数を単元とする。$R=k[z]/(A(z))$ に最初の変数を制限する。
$A(z)$ の項は消えるので、$B(z)$ の単元倍を除き $w$ 側の係数の包は
$1$ と $(1-t)zR$ を含む。実際 $A(w)/(w-z)$ は $R[w]$ で多項式となり、
その係数は最高次から三角で $R$ を張る。核に現れるのは
$A(w)+(1-t)z A(w)/(w-z)$ で、$w^m$ の係数からまず定数1が得られる。
$R/zR$ は高々一次元だから全 $m$ 次元が得られる。
尖点では $A(0)\ne0$ の側を選び、$1/z$ の平行移動で無限遠も根から外す。
二変数の $zw/(z-w)$ はこの移動で不変であり、二項式の自明化だけが変わる。
同じ商に現れる式は
$(1/t+z)A(w)+z^2 A(w)/(w-z)$ である。
$z$ が単元なので再び全 $R$ を張る。

従って階数1なら両切断が比例するか、共通三次因子と残余一次式を持つ。
後者で残余を $z-u,z-v$ とすると、残余二行行列式は節点で
$-(u-v)(t^2u-v)$、尖点で $(u-v)(2tuv-u+v)/t^2$ である。
同次化により無限遠も含む。比例しない場合の残余は従って
節点で $v=t^2u$、尖点で $1/v=1/u+2t$ となる。
比例する切断は節点で両端、尖点で0に二重に消える。
節点なら $u=v=0$、尖点なら $u=v=0$ に取って一次因子を一つ除けば
同じ表示ができる。従って以下の正規化入射は導手上の点を含め\emph{全ての台点}を覆う。
ここで非被約な底からの持上げは主張しない。
\end{proof}

\begin{proof}[定数スカラーを取り除く]
残余点の底を $\mathbb P^1$ 全体とする。節点では $[r:s]$ 上で
$a=(rz-s)g,\ b=\lambda(rz-t^2s)g$ とする。
共通三次式への条件は $r g_3=-t s g_0$ で両側とも同じである。
尖点では $a=(r-sz)g,\ b=\lambda(r-(s+2tr)z)g$ とし、
共通条件は $r g_1=(s+tr)g_0$ である。従って両側の三空間は底全体で同型な束で、
スカラー $\lambda$ の完備化は積 $\mathbb P^1\times\mathbb P^1$ である。
元の $a$ の三空間束 $A$ は四空間の部分束で次数 $-4$ である。
二つの同定した束から $k^4$ への単射を $T$ と第二区分の包含で作る。
各台点はこれらの差の鉛筆に核を与える。スカラー0と無限遠は単射である。

退化に曲線があれば $u$ または $[r:s]$ の底の上で水平であり、
スカラー射影は射影的で0と無限遠を避けるから一定である。
一定値 $\lambda_1$ の一般核から線部分束を取り、$A$ に飽和させる。
二写像はこの直線上で一定比で一致し、その像も部分束である。
両側をこの直線とその像で割ると再び二単射の鉛筆となる。
$\lambda\ne\lambda_1$ の各ファイバーでは元の核から商の核への写像は同型である。
左の直線成分の係数 $\lambda-\lambda_1$ が単元なので一意に戻せる。
これを残る退化が有限になるまで繰り返す。取り除いた回数を $j\le3$、
取り除いた部分束全体の次数の負を $b$ とする。これはもとの定数束の部分束であり $b\ge0$。
新しい底上の源と的の階数は $r'=3-j,r'+1$、次数は $b-4,b$ である。

$r'>0$ で退化が有限ならその長さは $j=0,1,2$ で順に
$16,12-b,8-b$ である。実際、二番目の射影直線の超平面類を $h$、最初のを $f$ とすれば、
源に $\mathcal O(-h)$ を掛けた写像の最大小行列式 の期待余次元は2で、
長さは $c_2(\mathrm{target}-\mathrm{source})$ である。
$f^2=h^2=0,\int fh=1$ と、源の第一類 $(b-4)f-r'h$、
第二類 $-(r'-1)(b-4)fh$ を代入すればよい。
これらの長さは取り除いたスカラーの平面外の台点数の上界にもなる。
核が一本なら一点である。核が二次元なら入力直線上に高々三点で、
退化イデアルは単元行列を消した後に全て二次以上ゆえ局所長は3以上である。
核が三次元の元の鉛筆では入力平面全体が階数1となり、整な $D$ に反する。
商での核の持上げは線型であり、取り除いた全スカラー外でこの議論が使える。
$j=0$ は既に二十に届かない。

一定スカラーの射影は入力平面に入る。節点では $b_0=t^2\lambda a_0$、尖点では
$b_0=\lambda a_0$ であり、$a_0,b_0$ は独立である。
従属なら共通の零平面で正規化の特異枝が共通根となり、結果式 $d$ の既約性に反する。
この平面条件は節点の $s=0$、尖点の $r=0$ でも成立する。
そこで $a_0=b_0=0$ が共通三次因子の貼合条件から出るためである。
一平面の台点数は高々7なので、$j=1$ は $7+12-b<20$、
$j=2$ で同じスカラーなら $7+8-b<20$ である。
$j=3$ は全ての台点が高々三平面に入り、三次評価の単射性に反する。
\end{proof}

\begin{proof}[二つの異なる固有直線]
$j=2$ で異なる二スカラーなら $b\ge3$ である。
実際、取り除いた階数2の部分束では二つの像束が一致し、
二つの相異なる定数固有値の射影子で元の鉛筆の二固有線束に分裂する。
各線束は入力空間の直線を与える非定数射である。
一点に固定されれば対応する非零四次式が全残余点 $u$ で消えてしまう。
二つの射の超平面次数の和が $b$ なので、両方が次数1となる場合だけを除けばよい。

次数1なら $a(z;r,s)$ の係数は $r,s$ の一次式であり、恒等的に指定の残余一次因子で割れる。
因子は $(r,s)$ に原始なので、その商 $g$ はこれらの変数に依らない一つの三次式である。
節点条件は $g_0=g_3=0$、尖点条件は $g_0=g_1=0$ を強制する。
従ってその入力直線は $a_0=b_0=0$ の共通直線全体となる。
両方が次数1なら同じ二空間 $\langle g,z g\rangle$ を与え、$g$ はこの二空間の
最大公約因子なので定数倍まで同じである。
節点では $r$ の係数、尖点では $s$ の係数を比べれば $T(zg)=\lambda zg$ であり、
同じ $g$ で二つの異なるスカラーは取れない。
従って平面上と外の台点を合わせても $2\cdot7+8-b\le19$ である。
これで全ての鉛筆を尽くした。
\end{proof}

\path{computations/verify_singular_residual_pencils.py} は二つの残余二行式の因子分解と
商の三つのChern数を記号計算で確認する。
